\chapter{1985年广东卷（文）}

\section{填空题}

\begin{problem}

\(\log_{3}5\cdot\log_{5}9 \) 的值等于\fillinblank{}.
\begin{answer}
\(2\)
\end{answer}

\begin{solution}
利用换底公式，
\(\log_{3}5\cdot\log_{5}9=\dfrac{\ln5}{\ln3}\cdot\dfrac{\ln9}{\ln5}=\dfrac{\ln9}{\ln3}=2\). 因此填 \(2\).
\end{solution}
\end{problem}

\begin{problem}

设数列 \(\{a_{n}\}\) 的首项 \(a_{1} = 8\)，且满足 \(a_{n} = 2a_{n-1}\)（\(n = 2\)，\(3\)，\(\cdots\)），则 \(a_{4}\) 的数值是\fillinblank{}.
\begin{answer}
\(64\)
\end{answer}

\begin{solution}
由递推关系逐项计算，\(a_{2}=2a_{1}=16\)，\(a_{3}=2a_{2}=32\)，\(a_{4}=2a_{3}=64\). 因此填 \(64\).
\end{solution}
\end{problem}

\begin{problem}

复数 \(\sqrt{3}-\i \) 的三角形式是\fillinblank{}.
\begin{answer}
\(2\left(\cos\frac{11\pi}{6}+\i\sin\frac{11\pi}{6}\right)\)
\end{answer}

\begin{solution}
复数 \(z=\sqrt3-\i\) 的模为 \(|z|=\sqrt{3+1}=2\). 它的实部为正、虚部为负，辐角在第四象限，且 \(\cos\theta=\frac{\sqrt3}{2}\)、\(\sin\theta=-\frac12\)，故可取 \(\theta=\frac{11\pi}{6}\). 所以
\(z=2\left(\cos\frac{11\pi}{6}+\i\sin\frac{11\pi}{6}\right)\)，与答案一致.
\end{solution}
\end{problem}

\begin{problem}

如果 \(f\left(\frac{1}{x}\right)=\frac{x}{1-x^{2}}\)，则 \(f(x)=\)\fillinblank{}.
\begin{answer}
\(\frac{x}{x^{2}-1}\)
\end{answer}

\begin{solution}
令 \(t=\frac1x\)，则 \(x=\frac1t\). 在原式有意义的范围内，
\(f(t)=\dfrac{1/t}{1-1/t^{2}}=\dfrac{1/t}{(t^{2}-1)/t^{2}}=\dfrac{t}{t^{2}-1}\). 将字母 \(t\) 改写为 \(x\)，得 \(f(x)=\dfrac{x}{x^{2}-1}\). 这里原条件要求代换时 \(x\ne0\)，\(\pm1\)，所得函数式在相应定义范围内成立.
\end{solution}
\end{problem}

\begin{problem}

如果一个圆锥底面的半径增大到原来的 3 倍，但要保持它的体积不变，则其高必须缩小到原来的\fillinblank{}.
\begin{answer}
\(\frac{1}{9}\)
\end{answer}

\begin{solution}
设圆锥原来的底面半径和高分别为 \(r\)、\(h\)，则体积为 \(V=\frac13\pi r^{2}h\). 半径变为 \(3r\) 后，若高为 \(h'\)，由体积不变得
\(\frac13\pi(3r)^{2}h'=\frac13\pi r^{2}h\)，所以 \(9h'=h\)，即 \(h'=\frac19h\). 因此填 \(\frac19\).
\end{solution}
\end{problem}

\begin{problem}

\(\cos \frac{\pi}{5} + \cos \frac{2\pi}{5} + \cos \frac{3\pi}{5} + \cos \frac{4\pi}{5}\) 的值等于\fillinblank{}.
\begin{answer}
\(0\)
\end{answer}

\begin{solution}
利用 \(\cos(\pi-\theta)=-\cos\theta\)，有 \(\cos\frac{4\pi}{5}=-\cos\frac{\pi}{5}\)，以及 \(\cos\frac{3\pi}{5}=-\cos\frac{2\pi}{5}\). 因而四项两两相消，原和为 \(0\).
\end{solution}
\end{problem}

\begin{problem}

\(\displaystyle\lim_{n\to\infty}\frac{1 + 4 + 7 + \cdots + (3n - 2)}{n(n + 1)}\) 的值等于\fillinblank{}.
\begin{answer}
\(\frac{3}{2}\)
\end{answer}

\begin{solution}
分子是首项为 \(1\)、公差为 \(3\)、项数为 \(n\) 的等差数列之和，因此
\(1+4+7+\cdots+(3n-2)=\frac{n[1+(3n-2)]}{2}=\frac{n(3n-1)}{2}\). 所以
\(\displaystyle\lim_{n\to\infty}\frac{1+4+7+\cdots+(3n-2)}{n(n+1)}=\lim_{n\to\infty}\frac{3n-1}{2(n+1)}=\frac32\).
\end{solution}
\end{problem}

\begin{problem}

在一个长方体中，和其中一条确定的棱相垂直的棱的总数是\fillinblank{} 条.
\begin{answer}
\(8\)
\end{answer}

\begin{solution}
长方体的棱只有三个方向，任取一条固定的棱后，与它同方向的另外三条棱均平行，不能计入；另外两个方向各有四条棱，并且这八条棱的方向都与固定棱垂直. 因此总数为 \(4+4=8\) 条.
\end{solution}
\end{problem}

\begin{problem}

双曲线 \(\frac{x^2}{3} - y^2 = 1\) 的两条渐近线所夹的锐角的大小是\fillinblank{} 弧度.
\begin{answer}
\(\frac{\pi}{3}\)
\end{answer}

\begin{solution}
双曲线可写为 \(\frac{x^{2}}{(\sqrt3)^{2}}-\frac{y^{2}}{1^{2}}=1\)，所以两条渐近线为 \(y=\frac{x}{\sqrt3}\) 和 \(y=-\frac{x}{\sqrt3}\). 设它们的锐角为 \(\theta\)，则
\(\tan\theta=\left|\dfrac{\frac1{\sqrt3}-(-\frac1{\sqrt3})}{1-\frac13}\right|=\sqrt3\). 因为 \(0<\theta<\frac\pi2\)，故 \(\theta=\frac\pi3\).
\end{solution}
\end{problem}

\begin{problem}

函数 \(y = \sin^{2}(2x) \) 的最小正周期是\fillinblank{}.
\begin{answer}
\(\frac{\pi}{2}\)
\end{answer}

\begin{solution}
由恒等式 \(\sin^{2}(2x)=\frac{1-\cos4x}{2}\)，函数的周期由 \(\cos4x\) 决定. 其最小正周期 \(T\) 满足 \(4T=2\pi\)，所以 \(T=\frac\pi2\).
\end{solution}
\end{problem}

\section{单选题}

\begin{problem}

函数 \(y = 2^{-x} \) 的图象的大致形状是（\quad）

\choicesfive
    {\choicebitmap[width=0.15\paperwidth]{img/1985/guangdong_liberal/q1_optA.png}}
    {\choicebitmap[width=0.15\paperwidth]{img/1985/guangdong_liberal/q1_optB.png}}
    {\choicebitmap[width=0.15\paperwidth]{img/1985/guangdong_liberal/q1_optC.png}}
    {\choicebitmap[width=0.15\paperwidth]{img/1985/guangdong_liberal/q1_optD.png}}
    {\choicebitmap[width=0.15\paperwidth]{img/1985/guangdong_liberal/q1_optE.png}}
\begin{answer}
E
\end{answer}

\begin{solution}
\(y=2^{-x}=\left(\frac12\right)^x\) 的函数值始终为正，且底数 \(\frac12\) 小于 \(1\)，所以图象从左向右递减；又有 \(f(0)=1\)，并且当 \(x\to+\infty\) 时 \(f(x)\to0\). 第五项的图象同时具有这些特征，故选第五项.
\end{solution}
\end{problem}

\begin{problem}

函数 \(y = \sqrt{\lg(1 - x)} \) 的定义域是（\quad）

\choicesfive
    {\((-\infty,+\infty) \)}
    {\((-\infty, 0] \)}
    {\((-\infty, 0) \)}
    {\((-\infty, 1] \)}
    {\((-\infty, 1) \)}
\begin{answer}
B
\end{answer}

\begin{solution}
要使 \(\sqrt{\lg(1-x)}\) 有意义，首先须有 \(1-x>0\)，并且根式内还须满足 \(\lg(1-x)\geqslant0\). 因为常用对数的底数大于 \(1\)，所以 \(1-x\geqslant1\)，即 \(x\leqslant0\). 定义域为 \((-\infty,0]\)，故选第二项.
\end{solution}
\end{problem}

\begin{problem}

以 \(A(1, 3)\)，\(B(-5, 1)\) 为端点的线段的垂直平分线的方程是（\quad）

\choicesfive
    {\(3x - y + 8 = 0\)}
    {\(3x + y + 8 = 0\)}
    {\(2x - y - 6 = 0\)}
    {\(2x + y + 2 = 0\)}
    {\(3x + y + 4 = 0\)}
\begin{answer}
E
\end{answer}

\begin{solution}
线段 \(AB\) 的中点为 \(M\left(-2,2\right)\)，斜率为 \(\dfrac{1-3}{-5-1}=\frac13\). 垂直平分线的斜率为 \(-3\)，所以其方程为
\(y-2=-3(x+2)\)，即 \(3x+y+4=0\). 因此选 E.
\end{solution}
\end{problem}

\begin{problem}

下列不等式中，成立的是（\quad）

\choicesfive
    {\(\sin \frac{3\pi}{8} < \sin \frac{7\pi}{8}\)}
    {\(\cos \frac{3\pi}{8} < \cos \frac{7\pi}{8}\)}
    {\(\csc \frac{3\pi}{8} < \csc \frac{7\pi}{8}\)}
    {\(\cot \frac{3\pi}{8} < \cot \frac{7\pi}{8}\)}
    {\(\tan \frac{3\pi}{8} < \tan \frac{7\pi}{8}\)}
\begin{answer}
C
\end{answer}

\begin{solution}
因为 \(\sin\frac{7\pi}{8}=\sin\frac\pi8\)，而 \(\frac{3\pi}{8}>\frac\pi8\) 且正弦函数在 \(\left(0,\frac\pi2\right)\) 上递增，所以 \(\sin\frac{3\pi}{8}>\sin\frac{7\pi}{8}\)，第一项错.
\(\cos\frac{3\pi}{8}>0>\cos\frac{7\pi}{8}\)，所以第二项错. 两个正弦值均为正，且 \(\sin\frac{3\pi}{8}>\sin\frac{7\pi}{8}\)，取倒数后不等号反向，得到 \(\csc\frac{3\pi}{8}<\csc\frac{7\pi}{8}\)，第三项对.
又 \(\cot\frac{3\pi}{8}>0>\cot\frac{7\pi}{8}\)，第四项错；\(\tan\frac{3\pi}{8}>0>\tan\frac{7\pi}{8}\)，第五项错. 故选第三项.
\end{solution}
\end{problem}

\begin{problem}

下列函数中，为奇函数的是（\quad）

\choicesfive
    {\(x^{2} + x^{3}\)}
    {\(\log_{x}x\)}
    {\(\sin x + 1\)}
    {\(\sqrt[3]{x} \)}
    {\(|x|^{3} \)}
\begin{answer}
D
\end{answer}

\begin{solution}
逐项检验奇函数条件 \(f(-x)=-f(x)\). \(x^{2}+x^{3}\) 的偶次项不满足该条件；\(\log_xx\) 的定义域为 \(x>0\) 且 \(x\ne1\)，不是关于原点对称；\(\sin x+1\) 含常数项；\(|x|^{3}\) 是偶函数. 对于 \(f(x)=\sqrt[3]{x}\)，定义域为 \(\R\)，且 \(f(-x)=\sqrt[3]{-x}=-\sqrt[3]{x}=-f(x)\)，所以它是奇函数. 故选第四项.
\end{solution}
\end{problem}

\begin{problem}

已知 \(180^{\circ} < \alpha < 270^{\circ}\)，且 \(\sin (270^{\circ} + \alpha) = \frac{4}{5}\)，则 \(\tan \frac{\alpha}{2} =\)（\quad）

\choicesfive
    {3}
    {2}
    {\(-2\)}
    {\(-3\)}
    {\(-\frac{1}{3}\)}
\begin{answer}
D
\end{answer}

\begin{solution}
由 \(\sin(270^{\circ}+\alpha)=-\cos\alpha\)，题设给出 \(\cos\alpha=-\frac45\). 因为 \(180^{\circ}<\alpha<270^{\circ}\)，故 \(\sin\alpha=-\frac35\). 半角公式给出
\(\tan\frac\alpha2=\dfrac{\sin\alpha}{1+\cos\alpha}=\dfrac{-3/5}{1-4/5}=-3\). 因此选第四项.
\end{solution}
\end{problem}

\begin{problem}

\(y = x^{2}\)（\(x \leqslant 0\)）的反函数是（\quad）

\choicesfive
    {\(y = \sqrt{x} \)}
    {\(y = \pm \sqrt{x} \)}
    {\(y = -\sqrt{x} \)}
    {\(y = \sqrt{-x} \)}
    {\(y = -\sqrt{-x} \)}
\begin{answer}
C
\end{answer}

\begin{solution}
原函数在 \(x\leqslant0\) 上为单调递减，值域为 \([0,+\infty)\). 由 \(y=x^{2}\) 且 \(x\leqslant0\)，解得 \(x=-\sqrt y\). 交换自变量与因变量，反函数为 \(y=-\sqrt x\)，其定义域为 \(x\geqslant0\). 故选第三项.
\end{solution}
\end{problem}

\begin{problem}

设 \(R = \{\text{实数}\}\)，\(M = \{\text{纯虚数}\}\)，\(C = \{\text{复数}\}\)，其中 \(C\) 是全集，则（\quad）

\choicesfive
    {\(R \cup \overline{R} = C\)}
    {\(\overline{M} \cup R = C\)}
    {\(M \cup R = C\)}
    {\(M \cap R = \{0\} \)}
    {\(C \cap \overline{R} = M\)}
\begin{answer}
A
\end{answer}

\begin{solution}
在全集 \(C\) 中，\(\overline R\) 是实数集 \(R\) 的补集，因此任意复数必属于 \(R\) 或 \(\overline R\)，且不可能同时属于二者，从而 \(R\cup\overline R=C\)，第一项正确. 对第二项，\(\overline M\cup R\) 不包含非零纯虚数；对第三项，\(M\cup R\) 不包含实部、虚部均非零的复数；对第四项，按本题“纯虚数”的定义，\(M\cap R=\varnothing\ne\{0\}\)；对第五项，\(C\cap\overline R=\overline R\) 是非实数复数集，不等于纯虚数集 \(M\). 因此只有第一项正确，故选第一项.
\end{solution}
\end{problem}

\begin{problem}

设向量 \(\overrightarrow{OZ}\) 对应于复数 \(-2\sqrt{3} + 4\i\). 把 \(\overrightarrow{OZ}\) 按顺时针方向旋转 \(60^\circ\) 得到 \(\overrightarrow{OZ_1}\)，则与向量 \(\overrightarrow{OZ_1}\) 对应的复数是（\quad）

\choicesfive
    {\(-3\sqrt{3} - \i\)}
    {\(\sqrt{3} + 5\i\)}
    {\(-2\sqrt{3} - 4\i\)}
    {\(2\sqrt{3} + 4\i\)}
    {\(1 + 3\sqrt{3}\i \)}
\begin{answer}
B
\end{answer}

\begin{solution}
顺时针旋转 \(60^{\circ}\) 等价于乘以复数
\(\cos(-60^{\circ})+\i\sin(-60^{\circ})=\frac12-\frac{\sqrt3}{2}\i\). 因而旋转后的复数为
\((-2\sqrt3+4\i)\left(\frac12-\frac{\sqrt3}{2}\i\right)=-\sqrt3+3\i+2\i+2\sqrt3=\sqrt3+5\i\). 所以选第二项.
\end{solution}
\end{problem}

\begin{problem}

某班上午要上语文，数学，体育，外语四门课，其中体育老师因故不能上第一节和第四节，则不同排课方案的种数是（\quad）

\choicesfive
    {10}
    {12}
    {20}
    {22}
    {24}
\begin{answer}
B
\end{answer}

\begin{solution}
体育课只能安排在第二节或第三节，有 \(2\) 种选择. 确定体育课的位置后，语文、数学、外语在其余三节中任意排列，有 \(3!=6\) 种. 因此排课方案数为 \(2\times3!=12\)，选 B.
\end{solution}
\end{problem}

\begin{problem}

\(\left(x + \frac{1}{\sqrt{x}}\right)^6\) 的展开式中常数项是（\quad）

\choicesfive
    {1}
    {6}
    {15}
    {20}
    {30}
\begin{answer}
C
\end{answer}

\begin{solution}
展开通项中，若取 \(k\) 个 \(x\)，取 \(6-k\) 个 \(x^{-1/2}\)，则该项的幂次为
\(k-\frac{6-k}{2}=\frac{3k-6}{2}\). 常数项要求幂次为零，所以 \(k=2\). 其系数为 \(\mathrm C_6^2=15\)，故选第三项.
\end{solution}
\end{problem}

\begin{problem}

不等式 \(|2x-3|>3 \) 的解集是（\quad）

\choicesfive
    {\(\{x\mid x>3\} \)}
    {\(\{x\mid x<0\} \)}
    {\(\{x\mid x>3 \text{且} x<0\}\)}
    {\(\{x\mid x>3\text{或}x<0\} \)}
    {\(\{x\mid 0<x<3\} \)}
\begin{answer}
D
\end{answer}

\begin{solution}
绝对值不等式等价于
\(2x-3>3\) 或 \(2x-3<-3\). 分别解得 \(x>3\) 或 \(x<0\)，所以解集为 \(\{x\mid x>3\text{或}x<0\}\)，选 D.
\end{solution}
\end{problem}

\begin{problem}

已知圆的方程为 \(x^{2} + y^{2} - 2x + 6y + 8 = 0\)，那么通过圆心的一条直线的方程是（\quad）

\choicesfive
    {\(2x - y + 1 = 0\)}
    {\(2x + y + 1 = 0\)}
    {\(2x - y - 1 = 0\)}
    {\(2x + y - 1 = 0\)}
    {\(y - 2x + 2 = 0\)}
\begin{answer}
B
\end{answer}

\begin{solution}
配方得 \(x^{2}-2x+y^{2}+6y+8=0\)，即 \((x-1)^{2}+(y+3)^{2}=2\). 所以圆心为 \((1,-3)\). 逐项代入圆心坐标，只有 \(2x+y+1=2-3+1=0\) 成立，故通过圆心的直线可取第二项.
\end{solution}
\end{problem}

\begin{problem}

已知椭圆的焦点是 \((1, 0)\) 和 \((-3, 0)\)，且短半轴长为 3，则椭圆的方程是（\quad）

\choicesfive
    {\(\frac{(x + 1)^2}{13} + \frac{y^2}{9} = 1\)}
    {\(\frac{x^2}{13} + \frac{y^2}{9} = 1\)}
    {\(\frac{(x - 1)^2}{13} + \frac{y^2}{9} = 1\)}
    {\(\frac{(x + 1)^2}{9} + \frac{y^2}{13} = 1\)}
    {\(\frac{x^2}{13} + \frac{(y - 1)^2}{9} = 1\)}
\begin{answer}
A
\end{answer}

\begin{solution}
两焦点的中点是椭圆中心 \((-1,0)\)，半焦距为 \(c=2\). 短半轴长为 \(b=3\)，故长半轴满足 \(a^{2}=b^{2}+c^{2}=9+4=13\). 长轴在 \(x\) 轴上，所以椭圆方程为
\(\dfrac{(x+1)^{2}}{13}+\dfrac{y^{2}}9=1\)，选 A.
\end{solution}
\end{problem}

\begin{problem}

在空间中，可以确定一个平面的条件是（\quad）

\choicesfive
    {两两相交的三条直线}
    {三条直线，其中的一条直线与另外两条直线分别相交}
    {三个点}
    {三条直线，它们两两相交，但不交于同一点}
    {两条直线}
\begin{answer}
D
\end{answer}

\begin{solution}
两条相交直线可以确定一个平面，但第一、第二、第五项都没有排除直线共点、平行或异面等情形，第三项的三个点也可能共线，均不能保证唯一确定平面.
若三条直线两两相交但不共点，任取其中两条直线确定一个平面；第三条直线分别与这两条直线相交于两个不同的点，因此第三条直线也在该平面内. 所以三条直线共面并确定一个平面，故选第四项.
\end{solution}
\end{problem}

\section{解答题}

\begin{problem}

证明 \(\frac{\tan(\alpha + \beta) - \tan\alpha}{1 + \tan\beta\tan(\alpha + \beta)} = \frac{\sin 2\beta}{2\cos^2\alpha}\).

\begin{solution}
在各式有意义且原分式分母不为零的范围内，利用正切加法公式，
\[
\begin{gathered}
\frac{\tan(\alpha + \beta) - \tan\alpha}{1 + \tan\beta\tan(\alpha + \beta)}
= \frac{\frac{\tan\alpha+\tan\beta}{1-\tan\alpha\tan\beta}-\tan\alpha}{1+\tan\beta\frac{\tan\alpha+\tan\beta}{1-\tan\alpha\tan\beta}}\\
= \frac{\tan\beta(1+\tan^{2}\alpha)}{1+\tan^{2}\beta}
= \frac{\frac{\sin\beta}{\cos\beta}\cdot\frac{1}{\cos^{2}\alpha}}{\frac{1}{\cos^{2}\beta}},\\
= \frac{\sin\beta\cos\beta}{\cos^{2}\alpha}
= \frac{\sin 2\beta}{2\cos^{2}\alpha}.
\end{gathered}
\]
\end{solution}
\end{problem}

\begin{problem}

用数学归纳法证明：\(1^{2}+3^{2}+5^{2}+\cdots+(2n-1)^{2}=\frac{1}{3}n(4n^{2}-1)\)（\(n\)为自然数）.

\begin{solution}
设命题 \(P(n)\) 为 \(1^{2}+3^{2}+5^{2}+\cdots+(2n-1)^{2}=\frac13n(4n^{2}-1)\). 当 \(n=1\) 时，左边为 \(1\)，右边为 \(\frac13(4-1)=1\)，所以 \(P(1)\) 成立.

假设 \(k\geqslant1\) 时 \(P(k)\) 成立，即
\(1^{2}+3^{2}+5^{2}+\cdots+(2k-1)^{2}=\frac13k(4k^{2}-1)\). 则当 \(n=k+1\) 时，
\[
1^{2}+3^{2}+5^{2}+\cdots+(2k-1)^{2}+(2k+1)^{2}
= \frac13k(4k^{2}-1)+(2k+1)^{2}
= \frac13k(2k+1)(2k-1)+(2k+1)^{2}
\]
于是
\[
\frac13k(2k+1)(2k-1)+(2k+1)^{2}
= \frac13(2k+1)(2k^{2}+5k+3)
= \frac13(2k+1)(k+1)(2k+3)
\]
再由
\[
\frac13(2k+1)(k+1)(2k+3)
= \frac13(k+1)\left[4(k+1)^{2}-1\right]
\]
所以 \(P(k+1)\) 也成立. 根据数学归纳法，原等式对所有自然数 \(n\) 都成立.
\end{solution}
\end{problem}

\begin{problem}

已知椭圆 \(C\) 和直线 \(l\) 的直角坐标方程分别为 \(C:x^{2}+2y^{2}-8=0\)，\(l:y=x-2\).

\begin{enumerate}
\item 证明椭圆 \(C\) 的一个焦点和短轴的一个端点落在直线 \(l\) 上.
\item 求直线 \(l\) 被椭圆 \(C\) 所截得线段的长度.
\item 如果动点 \(p\) 至坐标原点的距离与点 \(p\) 至直线 \(l\) 的距离相等，那么动点 \(p\) 的轨迹是什么曲线？请推导出它的方程.
\end{enumerate}

\begin{solution}
\begin{enumerate}
\item 【法一】将椭圆 \(C\) 的方程化为标准形式
\(\frac{x^{2}}{8}+\frac{y^{2}}{4}=1\).
得知椭圆的中心为坐标原点，短轴在 \(y\) 轴上，长半轴长 \(a=\sqrt{8}=2\sqrt{2}\)，短半轴长 \(b=\sqrt{4}=2\)，所以半焦距 \(c=\sqrt{a^{2}-b^{2}}=\sqrt{8-4}=2\)，从而两个焦点的坐标分别为 \((2,0)\) 和 \((-2,0)\)，又短轴两个端点的坐标分别为 \((0,2)\) 和 \((0,-2)\). 由直线 \(l\) 的方程 \(y=x-2\) 可知，当 \(x=2\) 时 \(y=0\)，当 \(x=0\) 时 \(y=-2\)，即椭圆 \(C\) 的焦点 \((2,0)\) 与短轴的端点 \((0,-2)\) 都在直线 \(l\) 上.

【法二】同【法一】知椭圆 \(C\) 的焦点为 \((2,0)\)，\((-2,0)\)，短轴的端点为 \((0,2)\)，\((0,-2)\). 由点 \((2,0)\) 与 \((0,-2)\) 所确定的直线的方程为
\(\frac{y-0}{-2-0}=\frac{x-2}{0-2}\)，
即
\(y=x-2\).
这就是直线 \(l\) 的方程，所以椭圆 \(C\) 的焦点 \((2,0)\) 与短轴的端点 \((0,-2)\) 都在直线 \(l\) 上.

\item 解方程组
\[
\begin{cases}
x^{2}+2y^{2}-8=0,\\
y=x-2.
\end{cases}
\]
消去 \(y\)，得
\(x^{2}+2(x-2)^{2}-8=0\)，
\(3x^{2}-8x=0\)，
解得
\(x_{1}=0\)，\(x_{2}=\frac{8}{3}\).
从而
\(y_{1}=-2\)，\(y_{2}=\frac{2}{3}\).
即直线 \(l\) 被椭圆 \(C\) 所截得的线段的端点坐标为 \((0,-2)\) 与 \(\left(\frac{8}{3},\frac{2}{3}\right)\)，该线段长为
\(\sqrt{\left(\frac{8}{3}-0\right)^{2}+\left(\frac{2}{3}+2\right)^{2}}=\frac{8\sqrt{2}}{3}\).

\item 由于动点 \(p\) 到坐标原点和到直线 \(l\) 的距离相等，可知动点 \(p\) 的轨迹是一条抛物线. 设动点 \(p\) 的坐标为 \((x,y)\)，依题意有
\(\sqrt{x^{2}+y^{2}}=\frac{|y-x+2|}{\sqrt{1^{2}+(-1)^{2}}}\).
等式两边均非负，平方不会引入或丢失轨迹上的点. 两边平方，得
\(x^{2}+y^{2}=\frac{1}{2}(y^{2}+x^{2}+4-2xy+4y-4x)\)，
即 \(x^{2}+y^{2}+2xy+4x-4y-4=0\). 令 \(u=x+y\)、\(v=x-y\)，方程化为 \(u^{2}+4v-4=0\)，即 \(u^{2}=-4(v-1)\)，所以轨迹确为抛物线，所求方程为 \(x^{2}+y^{2}+2xy+4x-4y-4=0\).
\end{enumerate}
\end{solution}
\end{problem}

\begin{problem}

已知直三棱柱 \(ABC-A_{1}B_{1}C_{1}\) 的侧棱 \(AA_{1}=4\mathrm{~cm}\)，它的底面 \(\triangle ABC\) 中有 \(AC=BC=2\mathrm{~cm}\)，\(\angle C=90^{\circ}\)，\(E\) 是 \(AB\) 的中点.

\begin{enumerate}
\item 求证：\(CE\perp AB_{1}\).
\item 求证：\(CE\) 和 \(AB_{1}\) 所在的异面直线的距离等于 \(\frac{2\sqrt{3}}{3}\mathrm{~cm}\).
\item 设截面 \(ACB_{1}\) 与侧面 \(ABB_{1}A_{1}\) 所成较小的二面角为 \(\alpha\)，试求 \(\cos\alpha\) 的值.
\end{enumerate}

\bitmapfigure[width=0.72\linewidth]{img/1985/guangdong_liberal/q4_fig1.png}
\begin{solution}
\begin{enumerate}
\item 【证法一】因为 \(ABC-A_{1}B_{1}C_{1}\) 是直三棱柱，因此 \(AA_{1}\perp\) 底面 \(ABC\). 因为 \(CE\subset\) 平面 \(ABC\)，因此 \(CE\perp AA_{1}\). 又 \(E\) 是等腰三角形 \(ABC\) 底边 \(AB\) 的中点，因此 \(CE\perp AB\). 过 \(E\) 作直线 \(m\parallel AA_{1}\)，则 \(m\) 在平面 \(ABB_{1}A_{1}\) 内，且 \(CE\perp m\). 直线 \(AB\) 与 \(m\) 相交于 \(E\)，所以 \(CE\perp\) 平面 \(ABB_{1}A_{1}\). 因为 \(AB_{1}\subset\) 平面 \(ABB_{1}A_{1}\)，因此 \(CE\perp AB_{1}\).

【法二】因为 \(ABC-A_{1}B_{1}C_{1}\) 是直三棱柱，因此侧面 \(ABB_{1}A_{1}\perp\) 底面 \(ABC\). 又 \(E\) 是等腰三角形 \(ABC\) 底边 \(AB\) 的中点，因此 \(CE\perp AB\)，故 \(CE\perp\) 平面 \(ABB_{1}A_{1}\). 因为 \(AB_{1}\subset\) 平面 \(ABB_{1}A_{1}\)，因此 \(CE\perp AB_{1}\).

【证法三】因为 \(E\) 是等腰三角形 \(ABC\) 底边 \(AB\) 的中点，因此 \(CE\perp AB\). 因为 \(ABC-A_{1}B_{1}C_{1}\) 是直三棱柱，因此 \(BB_{1}\perp\) 底面 \(ABC\). 又 \(AB\) 是 \(AB_{1}\) 在底面 \(ABC\) 上的射影，根据三垂线定理有 \(CE\perp AB_{1}\).

\item 【证法一】因为 \(CE\perp AB\)，\(CE\perp AB_{1}\)，因此 \(CE\perp\) 平面 \(ABB_{1}A_{1}\). 过 \(E\) 作 \(ED\perp AB_{1}\)，垂足为 \(D\)，\(AB\subset\) 平面 \(ABB_{1}A_{1}\)，因此 \(CE\perp ED\)，故 \(ED\) 是 \(CE\) 和 \(AB_{1}\) 所在的异面直线的距离. 在 \(\mathrm{Rt}\triangle ABB_{1}\) 和 \(\mathrm{Rt}\triangle ADE\) 中，\(\angle A\) 公共，因此 \(\triangle ABB_{1}\sim \triangle ADE\)，因此 \(\frac{ED}{BB_{1}}=\frac{AE}{AB_{1}}\)，即 \(ED=\frac{AE\cdot BB_{1}}{AB_{1}}\). 因为 \(AC=BC=2\)，\(\angle ACB=90^{\circ}\)，因此 \(AB=2\sqrt{2}\)，\(AE=\frac{1}{2}AB=\sqrt{2}\). 又在 \(\mathrm{Rt}\triangle ABB_{1}\) 中，\(BB_{1}=4\)，所以
\[
AB_{1}=\sqrt{AB^{2}+BB_{1}^{2}}=\sqrt{(2\sqrt{2})^{2}+4^{2}}=2\sqrt{6}
\]
因此
\(ED=\frac{\sqrt{2}\cdot4}{2\sqrt{6}}=\frac{2\sqrt{3}}{3}\mathrm{~cm}\).

【法二】同【证法一】至 \(ED\) 是 \(CE\) 和 \(AB_{1}\) 所在的异面直线的距离. 过 \(B\) 作 \(BF\perp AB_{1}\)，垂足为 \(F\)，则 \(ED\parallel BF\)，又 \(E\) 是 \(AB\) 的中点，因此 \(ED=\frac{1}{2}BF\). 在 \(\mathrm{Rt}\triangle ABB_{1}\) 中，\(BF\cdot AB_{1}=AB\cdot BB_{1}\)，即 \(BF=\frac{AB\cdot BB_{1}}{AB_{1}}\). 因为 \(AC=BC=2\)，\(\angle ACB=90^{\circ}\)，所以 \(AB=2\sqrt{2}\). 又 \(BB_{1}=4\)，\(AB_{1}=2\sqrt{6}\)，故
\[
ED=\frac{1}{2}BF=\frac{1}{2}\frac{2\sqrt{2}\cdot4}{2\sqrt{6}}=\frac{2\sqrt{3}}{3}\mathrm{~cm}
\]

\item 【法一】因为 \(ED\perp AB_{1}\)，\(CE\perp\) 侧面 \(ABB_{1}A_{1}\)，连结 \(CD\)，根据三垂线定理有 \(CD\perp AB_{1}\)，故 \(\angle CDE\) 是截面 \(CAB_{1}\) 和侧面 \(ABB_{1}A_{1}\) 所成的较小的二面角的平面角. \(CE\) 是 \(\mathrm{Rt}\triangle ABC\) 的斜边上的中线，故 \(CE=\frac{1}{2}AB=\sqrt{2}\). 又 \(ED=\frac{2\sqrt{3}}{3}\)，于是在 \(\mathrm{Rt}\triangle CED\) 中，
\[
\tan\angle CDE=\frac{CE}{ED}=\frac{\sqrt{2}}{2\sqrt{3}/3}=\frac{\sqrt{6}}{2}
\]
即 \(\tan\alpha=\frac{\sqrt{6}}{2}\). 又 \(\alpha\) 为锐角，
\[
\sec\alpha=\sqrt{1+\tan^{2}\alpha}=\sqrt{1+\frac{6}{4}}=\frac{\sqrt{10}}{2}
\]
\(\cos\alpha=\frac{1}{\sec\alpha}=\frac{\sqrt{10}}{5}\).

【法二】因为 \(CE\perp\) 侧面 \(ABB_{1}A_{1}\)，所以 \(\triangle CAB_{1}\) 在侧面上的射影是 \(\triangle EAB_{1}\). 设 \(\triangle CAB_{1}\) 的面积为 \(S_{1}\)，\(\triangle EAB_{1}\) 的面积为 \(S_{2}\)，则有 \(S_{2}=S_{1}\cos\alpha\). 在 \(\mathrm{Rt}\triangle CBB_{1}\) 中，\(B_{1}C=\sqrt{BB_{1}^{2}+BC^{2}}=\sqrt{4^{2}+2^{2}}=2\sqrt{5}\). 又 \(AB_{1}=2\sqrt{6}\)，\(AC=2\)，三角形 \(CAB_{1}\) 的半周长为 \(s=1+\sqrt{5}+\sqrt{6}\). 由海伦公式，\(S_{1}^{2}=s(s-2)(s-2\sqrt{5})(s-2\sqrt{6})=(10+2\sqrt{30})(-10+2\sqrt{30})=20\)，故 \(S_{1}=2\sqrt{5}\). 另一方面，
\[
S_{2}=\frac{1}{2}AB_{1}\cdot ED=\frac{1}{2}\cdot2\sqrt{6}\cdot\frac{2\sqrt{3}}{3}=2\sqrt{2}
\]
所以
\[
\cos\alpha=\frac{S_{2}}{S_{1}}=\frac{2\sqrt{2}}{2\sqrt{5}}=\frac{\sqrt{10}}{5}
\]
\end{enumerate}
\end{solution}

\end{problem}
